\documentclass[11pt,a4paper]{article}
\usepackage[margin=23mm]{geometry}
\usepackage{fontspec}
\setmainfont{DejaVu Serif}[ItalicFont={DejaVu Serif},ItalicFeatures={FakeSlant=0.18},BoldItalicFont={DejaVu Serif Bold},BoldItalicFeatures={FakeSlant=0.18}]
\setmonofont{DejaVu Sans Mono}[Scale=0.8]
\usepackage{babel}
\babelprovide[main,import]{ukrainian}
\usepackage{amsmath,booktabs,tabularx,longtable,microtype}
\usepackage[unicode,hidelinks]{hyperref}
\setlength{\emergencystretch}{3em}
\setlength{\parskip}{3pt}
\title{Індивідуальні контактні мережі:\\що вимірювали та що вдалося відтворити}
\author{Короткий пояснювальний огляд для роботи над SMPT}
\date{16 жовтня 2026 року}
\begin{document}
\maketitle
\noindent Мережі людей справді вимірюють і реконструюють. Найкраще перевірені приклади походять зі шкіл, лікарень та офісів. Для цілих міст частіше створюють синтетичних людей на основі перепису й мобільності. Ключове питання цього огляду: \textbf{що саме в реконструкції виявилося схожим на спостереження?} Однакові числа мешканців, схожа матриця контактів і схожа епідемічна крива --- різні рівні перевірки.

\section*{Що означають основні терміни}
В \textbf{індивідуальній мережі} вузол --- людина. Ребро може означати зафіксоване зближення, накопичену тривалість контакту, дружбу або лише змодельовану можливість зустрічі. Це потрібно уточнювати: дружба, фізичний контакт та поїздка не взаємозамінні.

\textbf{ASPL} --- середня кількість ребер у найкоротшому шляху між людьми. Для незв'язного графа потрібно вказати, чи рахували лише найбільшу компоненту. Це характеристика топології; вона не враховує час зустрічей. Шлях із послідовністю контактів у зворотному часовому порядку може існувати в статичному графі й бути неможливим для зараження.

\textbf{Кластеризація} показує, як часто знайомі однієї людини пов'язані між собою. Глобальна транзитивність --- частка замкнених трійок; середня локальна кластеризація --- середнє часток для окремих людей. Ці числа можуть суттєво різнитися. У порівнянні також мають збігатися часовий інтервал агрегації, поріг тривалості контакту та правила обліку ваг.

\section*{1. Виміряна шкільна мережа: Salathé, 2010}
\textbf{Дані.} Носимі датчики виміряли близькість між 788 учасниками протягом одного звичайного навчального дня; охоплення становило близько 94\%. У денному графі вузли --- люди, вага ребра --- сумарна тривалість їхньої близькості. Це виміряна мережа в межах конкретного протоколу датчиків~\cite{salathe}.

\textbf{Що отримали.} За різних порогів мінімальної тривалості контакту автори спостерігали короткі шляхи, значну кластеризацію й велику компоненту зв'язності. Одного універсального числа ASPL тут немає: граф змінюється після вилучення коротких контактів. Епідемічні симуляції на зваженому графі узгоджувалися зі шкільними даними про відсутність через гарячку під час H1N1. Це перевірка поведінки моделі, а не спостереження всіх актів передачі. Один день вимірювань не визначає довготривалу мережу будь-якої школи.

\section*{2. Синтетичне місто: Eubank, 2004}
\textbf{Побудова.} Для Portland використали перепис, землекористування та дані про пересування: близько 1.6~млн синтетичних людей, 181~тис. місць діяльності й 24-годинний розклад. Спочатку це мережа «люди--місця», з якої отримують контакти між людьми, що перебувають у тих самих місцях. Такі персональні ребра змодельовані, а не виміряні датчиками~\cite{eubank}.

\textbf{Що перевірили.} Автори описують персональний граф як подібний до малого світу й аналізують розподіл ступенів та кластеризацію. Перевірка TRANSIMS стосувалася транспортних потоків, заповнення місць, активностей та інших міських агрегатів. Це підтверджує правдоподібність мобільності, але не встановлює правильності кожної пари контактних осіб. Точних чисел ASPL чи глобальної кластеризації тут не наводимо: у перевіреному тексті для цього немає однозначного скалярного значення. Робота показує практичний шлях від міських даних до персонального графа.

\section*{3. Міські мережі: Guarino, 2021}
\textbf{Побудова.} Для Florence, Viterbo та Sabaudia поєднали просторовий розподіл населення, вік, домогосподарства й вікові матриці змішування. Людей з'єднували домашніми та дружніми зв'язками; імовірність дружби залежала від відстані, віку й індивідуальної товариськості~\cite{guarino}.

\textbf{Що отримали.} Проаналізовано ASPL, кластеризацію, ступені та компоненти. Нижче --- два приклади для синтетичної Florence з 363\,060 людьми. Це результати генератора, а не виміряні параметри реального графа міста. Домашні кліки збільшували кластеризацію; ребра не звіряли з повним реальним списком дружніх пар. Дружній зв'язок ще потребує окремої моделі фізичних контактів.

\begin{table}[htbp]
\centering\small
\begin{tabular}{rrrrrr}
\toprule
$\mu$ & $\langle k\rangle$ & ASPL & $C_{\mathrm{global}}$ & $\overline C_{\mathrm{local}}$ & Гігантська компонента\\
\midrule
5 & 6.82 & 6.72 & 0.053 & 0.074 & 98.8\%\\
10 & 11.81 & 5.36 & 0.017 & 0.027 & 99.9\%\\
\bottomrule
\end{tabular}
\caption{Вибрані рядки таблиці~6 у Guarino et al.: $\mu$ --- задане середнє число друзів; однакова товариськість, показник спадання ймовірності з відстанню~0.5. Середні за 10 побудовами. ASPL наведено за авторами; оскільки граф незв'язний, перед повторним використанням потрібно перевірити область його усереднення в реалізації.}
\end{table}

\section*{4. Реконструкція з матриць: Machens, 2013}
\textbf{Дані й реконструкція.} Виміряли контакти 119 осіб за тиждень у педіатричному відділенні. Спільноти --- професійні та соціальні ролі: пацієнти, лікарі, медсестри, помічники й супроводжувачі. З цих самих даних отримали матрицю середніх контактних тривалостей між ролями та багатшу матрицю розподілів. В останній кожна пара ролей задає розподіл персональних ваг, включно з нульовими контактами~\cite{machens}.

\textbf{Що збіглося.} Самі середні завищували розмір модельної епідемії й неправильно ранжували групи ризику. Матриця розподілів значно краще відтворювала розподіли розміру великих спалахів і ризики для ролей порівняно із симуляціями на детальній спостереженій мережі. Проте імовірність згасання не збігалася повністю з часовим еталоном. Це \textbf{подібність результатів симуляцій на різних представленнях одних контактних даних}; реальні спалахи тут не були еталоном. Також це реконструкція з контактних агрегатів, а не з транспортних OD-потоків.

\section*{5. Пропущені контакти: Génois, 2015}
\textbf{Постановка.} Використали спостережені контакти школи, офісу та конференції. Частину людей штучно вилучали рівномірно випадково, а потім додавали синтетичних людей і контакти. Найповніший варіант зберігав щільність зв'язків, групову структуру, неоднорідність ваг та статистику часу контактів~\cite{genois}.

\textbf{Що збіглося.} Розподіли модельних розмірів спалахів ставали значно ближчими до результатів на повніших спостережених даних, ніж при використанні самої неповної мережі. Відновлення лише щільності або середнього ступеня було недостатнім і могло завищувати ризик. Навіть найкращий варіант систематично завищував максимальні розміри спалахів і в більшості випадків --- імовірність та середній розмір великих спалахів: дрібні структури колег і друзів не відтворювалися. Автори не намагалися вгадати справжні пропущені пари. Результат залежить від припущення рівномірного відбору; «повний» датасет також не охоплював усю реальну популяцію.

\section*{6. Вікове змішування: Mistry, 2021}
\textbf{Побудова.} Із переписів та опитувань створили синтетичні популяції з домогосподарствами, школами й робочими середовищами; із них вивели вікові контактні матриці. Основне охоплення --- 277 субнаціональних регіонів дев'яти країн~\cite{mistry}.

\textbf{Що збіглося.} На незалежних від калібрування опитуваннях Pearson-кореляція між елементами нормованих синтетичних і спостережених матриць становила \textbf{0.92 для Франції, 0.90 для Японії та 0.80 для Shanghai}. Відтворювалися контакти однолітків і характерні вікові структури. Це кореляція агрегованих матриць, не точність відновлення індивідуальних ребер і не частка правильно вгаданих людей. Отримана відповідність обґрунтовує використання демографії для побудови структури змішування, але не встановлює правильність ASPL, кластеризації чи мостових осіб кожного міста.

\section*{7. Мобільність і зараження: Chang, 2021}
\textbf{Дані й модель.} Використали погодинні переміщення, що представляють близько 98~млн людей у десяти агломераціях США. Граф двочастковий: житлові переписні території та місця відвідування; ребро --- потік відвідувачів. Змішування відбувається через спільні місця з урахуванням їхньої площі, заповнення й тривалості відвідувань~\cite{chang}.

\textbf{Що збіглося.} Модель відтворювала зареєстровані криві випадків. Після калібрування до 14~квітня автори також перевірили відповідність випадків за відкладений період 15~квітня--9~травня 2020 року, використовуючи мобільність відповідного часу. Це важливий приклад епідемічної перевірки моделі на потоках. Але вузли тут --- \textbf{території та місця, а не окремі люди}: це метапопуляційна модель, яка не відновлює й не перевіряє список персональних контактних ребер. Вона показує, які дані про спільне перебування корисні понад прості числа поїздок між районами.

\section*{Що означає «схоже до реального»}
Порівняння з реальною чисельністю та віком перевіряє \textbf{популяцію}. Порівняння контактних матриць, ступенів, ASPL і кластеризації перевіряє певні \textbf{властивості мережі}. Подібність епідемій на реконструкції та спостереженому графі перевіряє \textbf{динамічну придатність реконструкції}. Відповідність зареєстрованим спалахам --- ще один, окремий рівень. Жоден із них автоматично не доводить інші.

Для нашого задуму найближчі Machens і Génois: вони прямо досліджують, що втрачається, коли від персональних контактів залишаються групові агрегати. Eubank і Chang показують два різні способи використання мобільності: синтез індивідуальних активностей і агреговане змішування через місця. Guarino та Mistry корисні для локальної структури й контактних обмежень.

\textbf{Лише потоків між спільнотами недостатньо для єдиного відновлення персонального графа.} Потрібні правила, хто саме їздить, де зустрічається й з ким контактує. Люди з міст $A$ і $B$ можуть контактувати на роботі в $C$, хоча прямого потоку $A\to B$ немає. ASPL і кластеризація додають обмеження, проте не визначають розміщення всіх міжспільнотних ребер. Для SMPT потрібно перевіряти також розподіл зовнішніх контактів між людьми та часи першого зараження спільнот на тому самому графі. Це наша методична інтерпретація огляду, а не результат уже виконаного нового експерименту.

\begin{thebibliography}{9}\footnotesize
\setlength{\itemsep}{1pt}\setlength{\parskip}{0pt}
\bibitem{salathe} Salathé M. et al. (2010). A high-resolution human contact network for infectious disease transmission. \emph{PNAS}, 107, 22020--22025. \href{https://doi.org/10.1073/pnas.1009094107}{DOI: 10.1073/pnas.1009094107}.
\bibitem{eubank} Eubank S. et al. (2004). Modelling disease outbreaks in realistic urban social networks. \emph{Nature}, 429, 180--184. \href{https://doi.org/10.1038/nature02541}{DOI: 10.1038/nature02541}.
\bibitem{guarino} Guarino S. et al. (2021). Inferring Urban Social Networks from Publicly Available Data. \emph{Future Internet}, 13, 108. \href{https://doi.org/10.3390/fi13050108}{DOI: 10.3390/fi13050108}.
\bibitem{machens} Machens A. et al. (2013). An infectious disease model on empirical networks of human contact: bridging the gap between dynamic network data and contact matrices. \emph{BMC Infectious Diseases}, 13, 185. \href{https://doi.org/10.1186/1471-2334-13-185}{DOI: 10.1186/1471-2334-13-185}.
\bibitem{genois} Génois M. et al. (2015). Compensating for population sampling in simulations of epidemic spread on temporal contact networks. \emph{Nature Communications}, 6, 8860. \href{https://doi.org/10.1038/ncomms9860}{DOI: 10.1038/ncomms9860}.
\bibitem{mistry} Mistry D. et al. (2021). Inferring high-resolution human mixing patterns for disease modeling. \emph{Nature Communications}, 12, 323. \href{https://doi.org/10.1038/s41467-020-20544-y}{DOI: 10.1038/s41467-020-20544-y}.
\bibitem{chang} Chang S. et al. (2021; online 2020). Mobility network models of COVID-19 explain inequities and inform reopening. \emph{Nature}, 589, 82--87. \href{https://doi.org/10.1038/s41586-020-2923-3}{DOI: 10.1038/s41586-020-2923-3}.
\end{thebibliography}
\end{document}
